% !TeX root = ../thuthesis-example.tex
\chapter{系统方案}
\section{总体设计}
本项目旨在打造一个功能完备、体验流畅且高度智能化的方言学习平台。系统的总体设计遵循以学习者为中心、以交互式练习为核心、以自适应规划为导向的原则。也就是说，平台围绕学习者的需求来设计功能，通过人机交互实现“学-练-评-荐”的闭环，不断根据用户表现动态调整教学策略，从而提供个性化指导，提升学习效果。

 平台必须提供从注册登录、选择方言目标，到知识呈现、发音练习、智能诊断、评估反馈、路径推荐等一系列完整的学习功能模块，确保学习闭环的各个环节都有对应支持。其中，知识呈现以可视化的知识节点网络为主要界面，引导用户按结构化的方式逐步学习方言知识；同时通过交互式知识图谱满足用户自由探索和溯源的需求。当用户进入具体的学习单元时，系统提供详细内容和标准发音示范音频，支持用户跟读练习和录音对比。
 
 平台设计突出人工智能对教学的深度支撑。借助本项目开发的智能诊断和自适应学习路径规划引擎，系统能够对用户录制的发音进行实时的声学分析与正确性判定，不仅给出对错结果，还能定位问题来源（如声母、韵母或声调错误），提供具体可行的改进建议。根据每次练习的结果，系统动态更新用户的学习状态，并生成专属的“错题本”，作为后续自适应学习路径规划的依据。整个设计保证每个学习者都能获得有针对性的反馈和推荐，实现真正的个性化学习。
 
 为保证交互流畅，系统必须确保语音诊断反馈在数秒内返回结果，即时响应用户操作。这要求后端架构具有高并发支持能力，即使多名学习者同时练习也能保持性能稳定。此外，发音诊断算法的准确性设计目标是达到专业水准，以赢得学习者信任。系统在可靠性方面要求高可用性，避免服务中断，并通过云端同步确保用户数据在不同终端上的一致性。易用性方面，界面设计力求简洁直观，针对不同年龄和技术背景的用户都能轻松使用；智能诊断给出的反馈措辞通俗易懂，尽量避免过度专业的术语。最后，整体架构在扩展性上要求良好：采用模块化设计，各功能组件彼此解耦，便于日后增加新的方言种类或升级优化某一功能而不影响其他部分。
\section{总体架构}
根据上述设计目标与要求，系统总体架构进行了分层和模块划分。整体上可以将平台架构分为三个层次：数据与知识层、核心服务层以及应用接口层。各层次结构如图\ref{fig:系统架构图}所示。
\begin{figure}[H]
    \centering
    \includegraphics[width=0.9\textwidth]{架构图.pdf}
    \caption{系统架构图}
    \label{fig:系统架构图}
\end{figure}
\subsection{数据与知识层}
该层是平台智能运作的基础，包括多模态方言声学知识图谱和现代汉语方言音档两部分。知识图谱涵盖了方言学习所需的所有核心知识点及其关系网络，包括词汇、音节、音系规则、发音示例及其多模态特征等。所有知识以结构化存储在高性能数据库中，支持高效的查询和推理。现代汉语方言音档则存储大量真实方言发音的音频和配套的文本语料，用于提供标准发音示范和语境例句支持。一方面，这一数据层通过知识检索接口供上层智能体查询，以提供权威的专业知识依据；另一方面，通过语料检索接口供发音比对和示例获取。

\subsection{核心算法层}
该层实现平台的智能功能和业务逻辑，是系统的中枢。核心服务层由若干模块化的服务组件构成，各组件通过标准接口协同工作，实现“学-练-评-荐”的闭环功能，其概览如图\ref{fig:核心}所示。

知识图谱查询与推理服务负责解析用户在学习过程中的查询请求（例如点击某知识点或提出问题）并在知识层检索相应的信息。该服务内置推理能力，可基于图谱关系回答复杂的语言学问题或找到相关的知识节点。当用户探索交互式知识图谱或使用智能问答助手时，该服务提供底层数据支持。

发音练习与诊断服务是平台最核心的服务之一，用于处理用户录音并给出智能诊断反馈。服务首先利用语音信号处理模块提取用户发音的声学特征（如利用Wav2Vec2模型获取隐式特征向量，或提取MFCC参数、基频曲线等显式特征），然后通过比对知识层中对应标准发音的特征来评估发音是否正确。比对过程结合机器学习模型和知识规则：例如利用预训练语音模型进行相似度向量检索，再通过精调的跨模态重排模型精确对齐用户发音与标准发音的差异。最终，诊断服务调用大型语言模型生成针对性的反馈报告，包括错误类型分析和改进建议。该报告可含多模态内容，如突出显示声调曲线差异、嵌入标准发音音频等，以提高反馈的直观性和可解释性。

学习者模型与自适应推荐服务负责根据用户实时更新的学习状态提供个性化的学习路径规划。该模块维护每个用户的学习者状态向量，记录其对各知识点的掌握度以及偏好特征。当用户完成练习后，服务会根据诊断结果更新其掌握度向量，并将错误记录到用户的错题本中。随后，服务通过算法分析用户在哪些知识点上存在薄弱环节，以及学习轨迹中可能缺失的环节，从知识图谱中筛选出最适合下一步学习的节点。推荐策略同时考虑知识的先后依赖关系和用户兴趣，提高学习路径的科学性和个性化。

文脉智能问答代理服务是一个面向高级查询的智能体协作模块，基于大型语言模型和链式智能体框架构建。当用户提出复杂的问题（例如比较不同方言中某词的差异，或追溯某方言现象的历史渊源），该服务会并行调用多个专职智能体对子任务进行处理，并检索知识图谱和方言档案库中的相关资料，最后汇总形成答案。智能问答代理能够结合文本、音频和图像等多模态证据回答问题，保证答复的专业性和可解释性，是对常规学习模式的有力补充。
\begin{figure}[H]
    \centering
    \includegraphics[scale=0.3]{核心层.pdf}
    \caption{核心算法层概览}
    \label{fig:核心}
\end{figure}
\subsection{应用接口层}
该层是用户直接交互的前端界面与应用程序部分，负责将核心服务层的功能以直观友好的方式呈现给用户。

多端客户端提供Web网页端为主的界面，兼顾移动App端以方便用户随时随地学习。界面设计遵循简洁一致的原则，突出方言学习的特色内容。例如，主学习界面以可视化知识网络展现当前方言的知识结构节点，用户点击节点即可查看内容并进入练习。知识图谱探索界面提供平滑缩放和拖拽漫游功能，用户可以像使用地图一样浏览整个方言知识版图，点击节点查看详细信息。练习界面提供录音控件与播放控件，录音结束后会自动显示系统诊断结果和反馈。所有界面元素与核心服务实时相连，实现所见即所得的交互体验。

前端通过富媒体组件呈现教学内容，包括文本说明、图片示例、音频播放和可视化图表等。例如，当诊断服务返回用户发音的声学分析时，界面能够动态绘制用户发音与标准发音的频谱图或基频曲线对比图，用户可据此直观了解问题所在。
应用层还包括错题本查看以及个性化推荐学习路径等界面。学习者可以随时查看自己的掌握情况，对系统推荐的学习路径给出反馈，这些交互会发送至后台的学习者模型服务，进一步完善用户状态模型，从而实现人机协同的自适应学习控制。

\section{技术路线}
为了达成本项目的建设目标，系统研发采用了循序渐进、层层深入的技术路线。从底层数据构建到上层应用开发，各阶段相互衔接又各有侧重，形成了一条清晰的开发主线，具体如图\ref{fig:技术}所示。本项目技术路线可以概括为“五个阶段、三个迭代”：

\begin{figure}[H]
    \centering
    \includegraphics[scale=0.3]{技术路线图.pdf}
    \caption{技术路线概览}
    \label{fig:技术}
\end{figure}
项目初期重点围绕方言知识的数字化和结构化展开。首先根据语言学理论设计出覆盖方言音系特点的本体Schema，包括词汇、音节、音系规则、发音人、多模态资产等核心实体类型及其声学关系。然后，通过权威方言辞典，音档、地方志和语言学专著等多源异构数据采集方言词汇、音系规则等结构化知识，并使用命名实体识别和关系抽取技术从文本中提取知识三元组。并行地，在专业录音环境下录制各目标方言的大量语音视频素材，结合使用Montreal Forced Aligner (MFA) 等工具将录音与文本强制对齐，从中提取音素级时间戳，对每个词条生成语音频谱、共振峰、基频曲线等特征。最后，将上述所得知识和特征数据融合入图数据库，完成多模态方言声学知识图谱的构建。这一阶段建立了系统的内容基石。

在具备高质量知识库后，中期研发围绕智能算法展开。本阶段首先开发文脉智能体框架 Dialectus，以大语言模型（例如GPT系列）为基础，结合Chain-of-Agents链式智能体方法，使系统能够利用知识图谱和方言语料库进行复杂问答和推理分析。与此同时，研发发音智能诊断算法：利用深度学习的语音表征模型（如Wav2Vec2）提取用户发音特征并建立向量索引，结合知识图谱中的标准发音进行相似度检索与比对，借鉴RAG思想让预训练语言模型参考检索结果来生成诊断反馈。随后，设计个性化自适应学习算法，通过引入学习者状态向量将用户的历史学习数据纳入决策，在CoT（Chain-of-Thought）提示和多步检索生成基础上，实现动态调整回答细粒度和教学顺序的能力。本阶段各算法先在离线环境以原型方式验证其有效性，包括构造模拟用户数据对算法进行调优（例如利用拒绝采样生成训练样本提升诊断生成模型质量），确保核心功能的准确性与稳定性。

在算法模块成熟后，进入集成阶段。根据总体架构设计，将各功能以微服务形式部署搭建。主要步骤包括：搭建图数据库服务器并导入最终知识图谱数据；部署语音特征向量检索服务；搭建大语言模型推理服务，并加载经过微调的诊断反馈生成模型；开发业务逻辑服务器，将用户管理、学习者状态更新、自适应推荐等逻辑封装为API。与此同时，前端团队采用合适的技术栈开发Web界面，实现知识网络可视化、录音交互控件、多媒体播放控件等，确保能够调用后端各项服务。通过一系列接口调试与数据格式对接，实现前后端各模块的无缝通信，打通从用户操作到后台处理再到结果呈现的全流程。

集成完成后，对平台进行全面的功能测试与用户体验评估。首先由研发人员依据测试用例对各模块进行验证，包括账户流程、知识检索正确性、发音诊断准确性、推荐路径合理性等。随后邀请小规模目标用户试用，收集反馈。在此基础上进入迭代优化：针对发现的问题调整算法参数或规则，例如优化诊断模型的提示词模板以提高反馈可读性，调整推荐策略以避免过难或过易的节点推荐等。性能方面通过压力测试发现系统瓶颈，进行针对性优化（如增加缓存、优化数据库索引）。同时利用真实用户互动数据触发“声学自进化”机制，不断丰富训练集并微调模型，从而提升算法效果。经过数轮迭代，平台功能和性能趋于稳定。

最后一步是在优化基础上将系统部署到生产环境，向目标用户群开放服务。在上线运行后，仍将持续监控系统表现。一方面通过自动化工具监测服务可用性和响应时间，确保平台稳定运行；另一方面继续收集用户学习数据和反馈，作为后续改进的依据。根据需要适时扩充知识图谱内容（如增添新方言或新词汇）以及替换更先进的AI模型，保持平台的行业领先性和生命力。通过这一系列循序推进且螺旋上升的技术路线，本项目最终成功将理论研究转化为实际应用，并实现系统性能与功能的持续演进。

\section{所用到的技术以及模块说明}
\subsection{多模态知识图谱构建与管理模块}
模块负责构建和维护方言声学知识图谱，是系统的知识中枢。本项目首先在技术上采用本体论+关系数据库的方式设计知识架构：借鉴W3C推荐的 OntoLex-Lemon 模型定义词汇、词形、词义等基本本体，将方言特有的音系和超音段特征扩展建模进去。通过定制的数据抽取脚本和NLP管道，从电子版方言词典和语料中自动提取实体和关系，解决了多音字、同义词等歧义问题。在多模态数据方面，模块运用了语音信号处理技术：利用 Montreal Forced Aligner 对语音和文本强制对齐，提取每个音节的起止时间；接着计算音频的梅尔频率倒谱系数 (MFCC)、共振峰(Formant)、基频(F0) 等声学特征。这些特征与原始音频、视频一起作为“多模态资产”关联到相应知识节点上。知识图谱数据最终存储在图数据库中，并通过Gremlin或Cypher查询语言提供检索接口。同时，采用知识图谱嵌入技术将实体和关系向量化，以便智能体模型高效利用知识。此模块的实现确保了平台拥有高质量、结构严谨且可机器理解的方言知识库。知识图谱的使用如\ref{fig:多模态知识图谱概览}所示
\begin{figure}[H]
    \centering
    \includegraphics[width=0.9\textwidth]{多模态知识图谱概览.pdf}
    \caption{多模态知识图谱概览}
    \label{fig:多模态知识图谱概览}
\end{figure}
\subsection{知识图谱可视化}
知识可视化与交互模块：此模块包含知识节点网络的可视化展示和交互式知识图谱浏览两大功能，主要技术难点在于图形渲染和人机交互。前端采用可视化库将知识图谱数据绘制成力导向图或环形节点图，使用渐进式渲染技术保证即使数百节点的图谱也能平滑显示和缩放。为提高交互体验，模块实现了碰撞检测和层次聚类算法，使得节点不会过度重叠且根据知识类别聚类分组，方便用户观察结构。交互方面，模块通过监听用户在图谱上的点击、悬停等操作，触发相应的知识查询请求，由后端知识图谱服务返回节点详情，再由前端以弹出卡片等形式动态展示信息。
\subsection{发音练习与智能诊断模块}
这是平台的核心功能模块，融合了语音处理与自然语言生成技术。首先在语音分析阶段，模块使用自监督预训练模型将用户录音转换为高维声学向量表示。同时，预先将知识图谱中所有标准发音音频经过相同模型编码，建立一个向量数据库。模块用户向量执行近似最近邻搜索，快速检索出最相似的几个标准发音候选。然后通过一个精细化的跨模态重排模型对候选进行相关度重新评分。该重排模型基于BERT等Transformer，将用户音频特征和候选标准发音特征拼接后输入，输出匹配分数，从中选出最匹配的标准发音作为对比基准，具体架构图如图\ref{fig:检索策略}。接下来，诊断模块提取用户发音与基准发音之间在关键声学参数上的差异，如能否对齐主要共振峰频率、基频曲线走向有无明显偏离等，并据此生成结构化的错误描述。最后，模块调用内置的讯飞星火执行检索增强生成，将结构化差异描述和知识图谱中相关规则文本一并输入，生成面向学习者的自然语言反馈。例如：“您的声调有误，标准发音应为高降调（参考声调曲线），建议练习发音时声音从高到低滑落”。生成过程中模型会适当引用知识图谱中的权威信息，确保解释的准确性和专业性。值得一提的是，该模块还能输出多模态反馈：通过调用图谱中的语图、口型视频等资产，结合Matplotlib等工具将其嵌入到反馈界面，使学习者得到全方位的指导。
\begin{figure}[H]
    \centering
    \includegraphics[width=0.9\textwidth]{tele.pdf}
    \caption{检索策略架构图}
    \label{fig:检索策略}
\end{figure}
\subsection{个性化评估与学习路径模块}
这是平台的个性化评估与学习路径规划模块，是实现自适应教学的核心组成部分。模块在设计上融合了学习者建模、路径优化、检索增强生成等技术，旨在为不同水平与偏好的用户提供高度契合的学习方案。

在学习者建模阶段，系统会为每位用户维护一个动态更新的学习者状态向量，其中包含知识掌握度向量与认知偏好向量两部分。掌握度向量的各个维度与知识图谱中的核心节点一一对应，用于表示用户对特定知识点（如声调类型、音韵规则）的熟练程度；偏好向量则反映用户在多模态信息呈现上的倾向，例如更偏好声谱图、口型视频还是音频示范。每次有效交互完成后，系统会根据练习成绩与交互行为，通过带有学习率的平滑更新公式调整掌握度，并根据用户与多媒体资产的交互频率更新偏好向量。

在学习路径规划阶段，系统会调用自适应检索生成模型，以多模态方言声学知识图谱与现代汉语方言音档作为唯一的知识源，根据当前状态向量生成多个可行的学习路径候选。候选路径会从三个维度进行量化评估：一是正确性，即路径是否能够引导用户获得符合事实的正确知识；二是适应性，即路径中内容的难度与呈现模态是否与用户当前状态匹配；三是效率，即路径在有限步骤内能否完成教学目标。系统通过启发式搜索或蒙特卡洛树搜索等方法，从候选路径中筛选出综合评分最高的方案。

最终输出为“个性化学习包”，由推理聚合智能体生成结构化的教学内容对象。学习包包含诊断摘要（指出本次练习中的关键问题及其声学依据）、核心教学模块（规则讲解、多模态示范、互动练习等）以及知识图谱链接。多模态示范部分可结合Matplotlib等可视化工具，将标准发音与用户发音的声学曲线对比、口型视频、语谱图等整合呈现，确保学习者从听觉与视觉双重角度获得指导。

在交付给前端渲染引擎后，学习包被转换为可交互界面呈现给用户。用户在使用过程中产生的各类操作（如播放音频、放大谱图、点击知识图谱节点、提交练习等）都会被实时记录并反馈给后端，以更新学习者状态向量。这一闭环机制确保了每一次交互都能推动状态模型演化，使下一轮教学规划更加精准，从而实现持续迭代优化的个性化学习体验。
\begin{figure}[H]
    \centering
    \includegraphics[width=0.9\textwidth]{个性化教学.pdf}
    \caption{多模态整合与闭环教学流程图}
    \label{fig:检索策略}
\end{figure}

\subsection{智能问答与知识检索模块}
此模块底层实现是文脉智能体 Dialectus 框架，其概览如图\ref{fig:文脉智能体}所示，具体来说由多个基于大型语言模型的子代理协同工作。模块首先对用户提出的问题进行意图解析和任务拆分（比如涉及音系分析的部分交由音韵智能体，涉及词汇用法的交给词汇智能体等）。各子代理分别向知识图谱或方言语料库执行检索，采用链式通信策略逐步汇总证据。过程中特别引入了奖励函数机制来引导检索路径：包括与查询相关性的奖励、新信息多样性的奖励和多模态资产价值的奖励，确保智能体既紧扣问题又发掘新颖有价值的证据。当所有子代理完成各自任务后，模块启动管理代理进行汇总推理，运用多模态思维链方法验证各证据间的一致性，形成最终答案。生成的答案可能包含文字解释、知识图谱节点引用，甚至音频或图像证据，使回答既详实又具说服力。该模块综合运用了知识图谱检索、向量语义检索、LLM Prompt Engineering等技术，在测试中表现出良好的专业准确性和上下文理解能力，对于复杂的方言问题能够给出高质量的回答。
\begin{figure}[H]
    \centering
    \includegraphics[width=0.9\textwidth]{文脉智能体框架.pdf}
    \caption{文脉智能体概览}
    \label{fig:文脉智能体}
\end{figure}
\subsection{自进化模型训练模块}
为了使系统的诊断反馈能力不断提升，本项目设计了后台自动运行的模型自适应进化模块。其实现基于CoRAG 和RAFT思想的融合，CoRAG概览如图\ref{fig:CoRAG}所示。当真实用户与系统交互并产生新的语音数据时，模块会自动生成训练样本：即将用户录音、系统检索到的标准发音及系统原反馈三者组成新的样本（三元组）。接着，模块利用预设的教师模型对这些样本的质量进行评估过滤。教师模型包括逻辑教师（由一个冻结的大型语言模型担任，审核反馈解释是否通顺合理）和声学教师（一个高精度声学对比模型，判断反馈与实际差错是否相符），只有同时通过这两方面判定的样本才被标记为高质量。这些优选样本被加入增量训练集，随后触发低秩微调 (LoRA)过程对诊断生成模型进行微调。LoRA技术使得模型微调高效且影响范围可控，从而支持系统以较高频率进行小幅参数更新。在这一模块持续作用下，系统的发音诊断模型能够“边学边教”，随着用户数量和时间的增长不断自我完善。据统计，自进化机制显著提高了模型对罕见发音错误的识别率，并优化了反馈措辞的针对性，充分体现出人工智能系统在数据驱动下自适应演化的优势。
\begin{figure}[H]
    \centering
    \includegraphics[width=0.9\textwidth]{coRAG.pdf}
    \caption{CoRAG概览}
    \label{fig:CoRAG}
\end{figure}
\section{所使用到的软件}
\paragraph{讯飞星火}
\begin{figure}[H]
    \centering
    
\includegraphics[scale=0.3]{xf.pdf}
    \caption{讯飞星火}
\end{figure}
讯飞星火是科大讯飞推出的新一代认知大模型平台，具备跨领域的语言理解、推理与生成能力。本项目将其作为智能体核心，结合方言声学知识图谱，实现对用户发音的多维度分析与个性化反馈。讯飞星火不仅在声学分析中提供高精度的语音分析，还能理解方言的声学特征，并将结构化的诊断结果转化为易懂的自然语言解释，形成闭环教学反馈链路。

\paragraph{腾讯云}
\begin{figure}[H]
    \centering
    
\includegraphics[scale=0.25]{tc.pdf}
    \caption{腾讯云}
\end{figure}
系统的云端服务全部部署在腾讯云平台，包括后端 API、知识图谱数据库、语音分析服务与多智能体调度模块。腾讯云提供了高可靠的计算、存储和网络资源，支持弹性伸缩和高并发访问，确保在大规模用户同时在线的情况下依然保持响应速度与稳定性。此外，腾讯云的安全机制与防护体系也为语音数据与用户隐私提供了保障。

\paragraph{Visual Studio Code}
\begin{figure}[H]
    \centering
    
\includegraphics[scale=0.3]{vscode.pdf}
    \caption{Visual Studio Code}
\end{figure}
前端开发团队采用 Visual Studio Code（VS Code）作为主要集成开发环境，配合多种插件完成 Web 与移动端界面的开发与调试。VS Code 的跨平台特性和丰富的生态插件（如实时预览、代码自动补全、版本控制集成）显著提升了开发效率。前端主要负责实现用户界面的交互逻辑、多模态反馈可视化以及与后端 API 的数据交互。
\paragraph{PyCharm}
\begin{figure}[H]
    \centering
    
\includegraphics[scale=0.3]{pycharm.pdf}
    \caption{pycharm}
\end{figure}
后端服务与算法模块的开发采用 PyCharm 作为主要 IDE，特别适合 Python 语言的项目管理与调试。PyCharm 提供了完善的调试工具、依赖管理和单元测试功能，帮助后端团队高效实现语音处理流水线、知识图谱检索、智能体任务调度及算法训练与推理功能。同时，PyCharm 与 Docker、Git 深度集成，便于代码版本管理与容器化部署。